% !TeX program = xelatex
% !TeX root = main.tex
\begin{frame}[plain]
  \vfill
  {\small\color{wine}\bfseries 科研组会\par}
  \vspace{5mm}
  {\fontsize{22}{26}\selectfont\bfseries On the Information Set of\\Individual Investors\par}
  \vspace{3mm}
  {\large\color{wine}个人投资者的信息集\par}
  \vspace{3mm}
  {\color{wine}\rule{.95\linewidth}{.75pt}}\par\medskip
  Toomas Laarits \quad Jeffrey Wurgler\par
  {\small New York University \hfill 26 年 4 月 Working Paper}\par
  \vspace{4mm}{\small 汇报人：杨博远\quad 27国金博0}\vfill
\end{frame}

\begin{frame}[label=authors]{作者与选文}
  \crumb{开场}{作者与版本}{NBER WP 33625 · SSRN 5187645}
  \leadrow{Jeffrey Wurgler}{NYU 教授；行为资产定价的资深学者。}
  \smallskip
  {\small
  \textbf{Baker--Wurgler (2006), \textit{Journal of Finance}}\par
  \textit{Investor Sentiment and the Cross-Section of Stock Returns}\par
  Baker--Wurgler；情绪指数。\par\smallskip
  \textbf{Barberis--Shleifer--Wurgler (2005), \textit{JFE}}\par
  \textit{Comovement}：SP500指数成分变动与股票收益联动。\par}
  \smallskip
  \leadrow{Toomas Laarits}{NYU Stern，助理教授。}
  \smallskip\small
  \begin{tabularx}{\linewidth}{@{}lY@{}}
    2025 Working Paper & \textit{The Research Behavior of Individual Investors}\\
    2026-04-27 版本 & \textit{On the Information Set of Individual Investors}
  \end{tabularx}
\end{frame}

\section{研究问题}
\sectionlead{section-1}{一、研究问题}{资产定价需要一个信息集。个人投资者实际获得的信息怎样观测？}

\subsection{1.1 信息集}
\begin{frame}[label=chain]{资产定价中的信息如何进入决策}
  \crumb{Part 1}{1.1 信息集}{引言 pp.\,2--3}
  \centering
  \begin{tikzpicture}[node distance=3mm]
    \node[storybox,text width=1.65cm] (menu) {信息供给};
    \node[storyfocus,text width=1.65cm,right=of menu] (see) {信息获取};
    \node[storybox,text width=1.65cm,right=of see] (belief) {信念形成};
    \node[storybox,text width=1.65cm,right=of belief] (order) {投资决策};
    \node[storybox,text width=1.65cm,right=of order] (price) {市场价格};
    \draw[storyarrow] (menu)--(see);
    \draw[grapharrow] (see)--(belief);
    \draw[storyarrow] (belief)--(order);
    \draw[storyarrow] (order)--(price);
  \end{tikzpicture}\par\raggedright
  \[
    P_{j,t}=\mathbb{E}\!\left[M_{t+1}X_{j,t+1}\mid\mathcal I_t\right].
  \]
  {\small 定价式，\(M\) 为随机折现因子，\(X\) 为支付，\(\mathcal I_t\) 为条件信息集。}
  \medskip
  \pair{\small\textbf{甲：}\; 查盈利与基本面\\
       \(\to\) 上调现金流预期 \(\to\) 买入}
       {\small\textbf{乙：}\; 看近期价格下行\\
       \(\to\) 预期跌势延续 \(\to\) 卖出}
  \finding{甲乙的期望价格不同，可能源于所依据的信息集不同。}
\end{frame}

\begin{frame}[label=literature]{资产定价模型中的信息集假设}
  \crumb{Part 1}{1.1 信息集}{引言与文献 pp.\,2--7}
  \small
  \begin{tabularx}{\linewidth}{@{}L{.31\linewidth}YY@{}}
    \toprule 文献 & 关注的信息假设 & 常见经验观察\\\midrule
    Markowitz／Sharpe & 预期收益和风险 M-V & 组合权重、收益与风险\\
    Hong--Stein (1999) & 新闻逐步扩散，引导动量交易 & 价格延续与反转\\
    Barberis et al. (2015) & 历史收益进入外推信念 & 预期调查与历史收益\\
    Barber--Odean (2008) & 显著事件和极端表现的影响 & 新闻、成交量与买单\\
    \bottomrule
  \end{tabularx}
  \finding{特定资产定价模型依赖特定的信息集与信息假设。}
\end{frame}

\subsection{1.2 注意力证据}
\begin{frame}[label=proxies]{注意力代理能看到什么}
  \crumb{Part 1}{1.2 注意力证据}{引言 pp.\,3--7}
  \small
  现有注意力的信息集测度，间接且零散。
  \medskip
  \begin{tabularx}{\linewidth}{@{}L{.36\linewidth}Y@{}}
    \toprule 数据 & 可见对象\\\midrule
    搜索量、新闻阅读量 & 受关注的时间和强度\\
    单券商点击与账户活动 & 站内查阅客户交易\\
    \mbox{\footnotesize\bfseries 本文的网址链接URL数据} & 同一投资者的跨站访问、研究主题与订单时点\\\bottomrule
  \end{tabularx}
  \medskip
  {\footnotesize 链接示例\enspace
    \texttt{http://finance.yahoo.com/q/ao?s=AAPL}}\par
  \finding{网址链接数据为投资者微观信息获取行为提供高质量代理。}
\end{frame}

\begin{frame}[label=dimensions]{信息集：用四个维度刻画}
  \crumb{Part 1}{1.2 注意力证据}{引言 pp.\,3--5}
  \centering
  \begin{tikzpicture}
    \node[storyfocus,text width=5.1cm] (root) at (0,0) {拆解同一投资者（家庭）的网页链接记录};
    \foreach \x/\branchname/\branchtext in {
      -4.08/{深度 (Depth)}/{研究多久？},-1.36/{主题 (Topics)}/{看哪些信息？},
       1.36/{异质性}/{股票和投资者如何不同？},4.08/{时机 (Timing)}/{何时、依何顺序？}} {
      \node[storybox,text width=2.28cm,minimum height=1.44cm] (b) at (\x,-2.08)
        {\textbf{\branchname}\\[-1pt]\scriptsize \branchtext};
      \draw[storyarrow] (root.south)--(0,-.91)-| (b.north);
    }
  \end{tikzpicture}\par\raggedright
  \finding{本文侧重测度与描述性实证，用多个维度刻画信息集。}
\end{frame}

\begin{frame}[label=agenda]{目录}
  \crumb{报告路线}{研究问题到贡献与评价}{}
  \begin{columns}[T,onlytextwidth]
  \column{.51\textwidth}
    \textbf{\jump{section-1}{一、研究问题}}\par
    {\small 1.1 信息集\\1.2 注意力证据}\par\medskip
    \textbf{\jump{section-2}{二、数据与测度}}\par
    {\small 2.1 浏览样本与 URL\\2.2 记录整理与交易归属}\par\medskip
    \textbf{\jump{section-3}{三、描述性统计}}\par
    {\small 3.1 研究投入\\3.2 主题与窗口}
  \column{.46\textwidth}
    \textbf{\jump{section-4}{四、实证比较}}\par
    {\small 4.1 股票与投资者\\4.2 交易时点与路径}\par\medskip
    \textbf{\jump{section-5}{五、贡献与评价}}\par
    {\small 5.1 证据与局限\\5.2 研究评价}
  \end{columns}
\end{frame}
